% Options for packages loaded elsewhere
\PassOptionsToPackage{unicode}{hyperref}
\PassOptionsToPackage{hyphens}{url}
%
\documentclass[
]{article}
\usepackage{amsmath,amssymb}
\usepackage{iftex}
\ifPDFTeX
  \usepackage[T1]{fontenc}
  \usepackage[utf8]{inputenc}
  \usepackage{textcomp} % provide euro and other symbols
\else % if luatex or xetex
  \usepackage{unicode-math} % this also loads fontspec
  \defaultfontfeatures{Scale=MatchLowercase}
  \defaultfontfeatures[\rmfamily]{Ligatures=TeX,Scale=1}
\fi
\usepackage{lmodern}
\ifPDFTeX\else
  % xetex/luatex font selection
\fi
% Use upquote if available, for straight quotes in verbatim environments
\IfFileExists{upquote.sty}{\usepackage{upquote}}{}
\IfFileExists{microtype.sty}{% use microtype if available
  \usepackage[]{microtype}
  \UseMicrotypeSet[protrusion]{basicmath} % disable protrusion for tt fonts
}{}
\makeatletter
\@ifundefined{KOMAClassName}{% if non-KOMA class
  \IfFileExists{parskip.sty}{%
    \usepackage{parskip}
  }{% else
    \setlength{\parindent}{0pt}
    \setlength{\parskip}{6pt plus 2pt minus 1pt}}
}{% if KOMA class
  \KOMAoptions{parskip=half}}
\makeatother
\usepackage{xcolor}
\usepackage{longtable,booktabs,array}
\usepackage{calc} % for calculating minipage widths
% Correct order of tables after \paragraph or \subparagraph
\usepackage{etoolbox}
\makeatletter
\patchcmd\longtable{\par}{\if@noskipsec\mbox{}\fi\par}{}{}
\makeatother
% Allow footnotes in longtable head/foot
\IfFileExists{footnotehyper.sty}{\usepackage{footnotehyper}}{\usepackage{footnote}}
\makesavenoteenv{longtable}
\setlength{\emergencystretch}{3em} % prevent overfull lines
\providecommand{\tightlist}{%
  \setlength{\itemsep}{0pt}\setlength{\parskip}{0pt}}
\setcounter{secnumdepth}{-\maxdimen} % remove section numbering
\ifLuaTeX
  \usepackage{selnolig}  % disable illegal ligatures
\fi
\usepackage[]{natbib}
\bibliographystyle{plainnat}
\IfFileExists{bookmark.sty}{\usepackage{bookmark}}{\usepackage{hyperref}}
\IfFileExists{xurl.sty}{\usepackage{xurl}}{} % add URL line breaks if available
\urlstyle{same}
\hypersetup{
  hidelinks,
  pdfcreator={LaTeX via pandoc}}

\author{}
\date{}

\begin{document}

\hypertarget{when-does-dynamic-weight-decay-help-a-unified-framework-analysis}{%
\section{When Does Dynamic Weight Decay Help? A Unified Framework
Analysis}\label{when-does-dynamic-weight-decay-help-a-unified-framework-analysis}}

\hypertarget{abstract}{%
\subsection{Abstract}\label{abstract}}

Weight decay is universally applied in deep learning optimization, yet
practitioners face a bewildering proliferation of dynamic weight decay
strategies---Cautious Weight Decay (CWD), Scheduled Weight Decay (SWD),
cosine schedules, AlphaDecay, ADANA, AdamO---with no systematic
understanding of \emph{when} these methods provide genuine benefit. We
introduce the \textbf{Phi Modulator Framework}, a unified mathematical
abstraction \(\varphi(t, \boldsymbol{\theta}, \mathbf{s})\) that
recovers all major dynamic weight decay methods as special cases along
four modulation axes: temporal, directional, spatial, and target-norm.
Through this lens, we conduct the first controlled comparison of seven
weight decay strategies under identical optimization conditions across
87+ experiments. Our central finding is a \textbf{conditional
indistinguishability result}: under standard AdamW settings
(\(\rho = \lambda/\eta \approx 0.5\)), all seven strategies---including
the degenerate case of no weight decay---are statistically
indistinguishable (accuracy range \(< 0.3\%\), all pairwise
\(p_{\text{adj}} > 0.05\) after Holm correction, \(n = 3\)). In stark
contrast, SGD exhibits significant sensitivity to weight decay presence
(constant vs.~no\_wd: Cohen's \(d = 10.29\),
\(p_{\text{adj}} = 0.002\)), yielding an approximately \(18.3\times\)
effect-size ratio relative to AdamW (Bootstrap BCa 95\% CI:
\([12.1\times, 28.7\times]\)); this ratio reflects both the optimizer
mechanism difference and the 100\(\times\) difference in \(\rho\) values
between standard AdamW (\(\rho = 0.5\)) and SGD (\(\rho = 0.005\))
configurations. We propose the \textbf{Phi Invariance Conjecture}, which
posits \(\rho = \lambda/\eta\) as the order parameter governing a regime
boundary: when \(\rho \lesssim 1\), AdamW's implicit \(\ell_\infty\)
constraint is hypothesized to render all modulation strategies
indistinguishable; as \(\rho\) increases, strategy choice becomes
progressively important. We further propose three standardized
diagnostic metrics---BEM, CSI, and AIS---as characterization tools for
describing operational differences between weight decay strategies. All
findings are currently validated on CIFAR-scale benchmarks with
batch-normalized architectures; large-scale and BN-free validation
remains an important next step.

\begin{center}\rule{0.5\linewidth}{0.5pt}\end{center}

\hypertarget{introduction}{%
\subsection{1. Introduction}\label{introduction}}

\hypertarget{motivation}{%
\subsubsection{1.1 Motivation}\label{motivation}}

Weight decay is among the most universally applied techniques in deep
learning optimization. Yet despite its ubiquity, the community lacks a
principled framework for choosing \emph{how} to apply weight decay over
the course of training.

The classical understanding treats weight decay as explicit L2
regularization (Krogh \& Hertz, 1991). This view has been progressively
undermined: Loshchilov \& Hutter (2019) demonstrated that L2
regularization and decoupled weight decay are not equivalent in adaptive
optimizers, leading to AdamW. D'Angelo et al.~(2024) showed weight decay
is never useful as explicit regularization in modern BN-equipped
architectures; instead, it acts as a training dynamics modifier. Xie \&
Li (2024) revealed that AdamW implicitly performs \(\ell_\infty\)-norm
constrained optimization. Kosson et al.~(2023) showed weight decay
induces a rotational equilibrium across layers.

These understandings have catalyzed a surge of dynamic weight decay
methods: SWD (Xie et al., 2023), CWD (Chen et al., 2026a), AdamWN
(Loshchilov, 2023), AlphaDecay (He et al., 2025), ADANA (Ferbach et al.,
2026), AdamO (Chen et al., 2026b). A critical problem pervades this
literature: \textbf{each method is evaluated in isolation}, using
different architectures, datasets, optimizers, and hyperparameter
protocols. No two papers share the same experimental conditions.

\hypertarget{research-gap}{%
\subsubsection{1.2 Research Gap}\label{research-gap}}

Answering \emph{when} dynamic weight decay helps is currently impossible
due to four gaps: no unified mathematical framework, no standardized
evaluation metrics, no controlled systematic comparison, and no theory
for when modulation becomes irrelevant.

\hypertarget{our-approach-the-rho-lambdaeta-lens}{%
\subsubsection{\texorpdfstring{1.3 Our Approach: The
\(\rho = \lambda/\eta\)
Lens}{1.3 Our Approach: The \textbackslash rho = \textbackslash lambda/\textbackslash eta Lens}}\label{our-approach-the-rho-lambdaeta-lens}}

The quantity \(\rho = \lambda/\eta\) appears implicitly in recent
theoretical analyses (Xie \& Li, 2024; Defazio, 2025; Wang \& Aitchison,
2024). Our contribution is to operationalize it as an \emph{empirically
testable regime boundary parameter}. The ratio \(\rho\) measures the
relative magnitude of the weight decay step compared to the gradient
step.

Our core finding: at AdamW settings (AdamW:
\(\lambda = 5 \times 10^{-4}\), \(\eta = 10^{-3}\), \(\rho = 0.5\); SGD:
\(\lambda = 5 \times 10^{-3}\), \(\eta = 0.1\), \(\rho = 0.05\)), seven
weight decay strategies spanning all four modulation axes produce
statistically indistinguishable results under AdamW (see Figure 1 for
the accuracy distributions and Figure 2 for the AdamW/SGD contrast). Yet
under SGD, weight decay presence produces a massive effect (Cohen's
\(d > 10\), \(p < 0.003\)). We attribute this asymmetry to AdamW's
implicit \(\ell_\infty\) constraint (Xie \& Li, 2024), which we
hypothesize absorbs weight decay perturbations in the low-\(\rho\)
regime; Section 7.2 discusses the interplay with batch normalization's
scale-invariance (D'Angelo et al., 2024).

The \(\rho\)-based perspective connects \(\tau^* = 1/\rho\) (Xie \& Li,
2024), \(R_* = \rho\) (Defazio, 2025), and EMA timescale
\(\approx 1/\rho\) (Wang \& Aitchison, 2024) through a unified dual
characterization (Remark 1, Section 4).

\hypertarget{contributions}{%
\subsubsection{1.4 Contributions}\label{contributions}}

We make four contributions, ordered from enabling to dependent:

\begin{enumerate}
\def\labelenumi{\arabic{enumi}.}
\item
  \textbf{Phi Modulator Framework.} A unified mathematical interface
  \(\varphi(t, \boldsymbol{\theta}, \mathbf{s})\) recovering all major
  dynamic weight decay methods as special cases along four modulation
  axes, enabling controlled comparison under identical optimization
  conditions.
\item
  \textbf{Empirical discovery: conditional indistinguishability.} Across
  87+ controlled experiments (7 methods \(\times\) 3 seeds \(\times\) 2
  datasets \(\times\) 2 optimizers), we detect no significant
  differences among weight decay strategies under standard AdamW
  settings (all pairwise \(p_{\text{adj}} > 0.05\), Holm correction,
  \(n = 3\)). SGD exhibits significant sensitivity (\(18.3\times\)
  effect-size ratio, Bootstrap BCa 95\% CI:
  \([12.1\times, 28.7\times]\); see Table 2). Results apply to
  batch-normalized architectures at CIFAR scale (ResNet-20, 270K
  parameters; large-scale validation addressed in Section 7.3).
\item
  \textbf{Phi Invariance Conjecture.} A trichotomy positing \(\rho\) as
  the regime boundary parameter: Regime I (\(\rho \lesssim 1\),
  strategy-invariant), Regime II (transitional), and Regime III
  (\(\rho \gg 1\), strategy-sensitive). Currently tested only at
  \(\rho = 0.5\); Regimes II and III remain empirically unverified.
\item
  \textbf{Diagnostic tools: BEM/CSI/AIS.} Standardized metrics providing
  the first vocabulary for characterizing operational differences
  between weight decay strategies. These metrics do not predict which
  method is best; they quantify budget deviations, coupling stability,
  and alignment informativeness.
\end{enumerate}

\hypertarget{paper-roadmap}{%
\subsubsection{1.5 Paper Roadmap}\label{paper-roadmap}}

Section 2 surveys related work. Section 3 introduces the Phi Modulator
Framework. Section 4 presents theoretical analysis. Section 5 describes
the experimental setup. Section 6 presents results including Figures
1--3. Section 7 provides discussion and limitations. Section 8
concludes.

\begin{center}\rule{0.5\linewidth}{0.5pt}\end{center}

\hypertarget{related-work}{%
\subsection{2. Related Work}\label{related-work}}

\hypertarget{weight-decay-as-a-dynamics-modifier}{%
\subsubsection{2.1 Weight Decay as a Dynamics
Modifier}\label{weight-decay-as-a-dynamics-modifier}}

Loshchilov \& Hutter (2019) demonstrated that L2 regularization and
decoupled weight decay produce fundamentally different behaviors in
adaptive optimizers, leading to AdamW. D'Angelo et al.~(2024) showed
weight decay is \emph{never} useful as explicit regularization in modern
BN-equipped architectures. Kosson et al.~(2023) showed weight decay
induces a \emph{rotational equilibrium}. Xie \& Li (2024) proved AdamW
implicitly performs \(\ell_\infty\)-norm constrained optimization with
constraint radius \(\tau^* = \eta/\lambda = 1/\rho\).

Defazio (2025) derived the steady-state gradient-to-weight ratio
\(R_* \approx \lambda/\eta = \rho\) for normalized layers---the same
quantity as Xie \& Li's \(\tau^*\) viewed from a gradient equilibrium
perspective. Although Defazio does not make this connection explicit,
both are dual characterizations of \(\rho\) (Remark 1, Section 4).
Defazio's focus is end-of-training instability correction; our novelty
is the Trichotomy regime structure and cross-optimizer comparison, not
the \(\rho\) quantity itself.

\textbf{A critical open question}: whether BN scale-invariance or
AdamW's \(\ell_\infty\) constraint is the primary driver of the observed
strategy invariance. The NoBN ablation is the highest-priority blocking
experiment (Section 7.3, Limitation 2).

\hypertarget{dynamic-weight-decay-methods}{%
\subsubsection{2.2 Dynamic Weight Decay
Methods}\label{dynamic-weight-decay-methods}}

We organize existing methods along the four Phi modulation axes.
\textbf{Temporal scheduling}: SWD/AdamS (Xie et al., 2023); ADANA
(Ferbach et al., 2026); alignment-dependent bounds (Sun et al., CVPR
2025). \textbf{Directional modulation}: CWD (Chen et al., 2026a); AdamO
(Chen et al., 2026b); SPD for fine-tuning (Tian et al., 2024)---outside
our scope. \textbf{Spatial modulation}: AlphaDecay (He et al., 2025).
\textbf{Target-norm control}: AdamWN (Loshchilov, 2023).

Our key distinction: we provide the first systematic evaluation of
\emph{when} existing methods provide genuine benefit under identical
conditions.

\hypertarget{evaluation-fragmentation}{%
\subsubsection{2.3 Evaluation
Fragmentation}\label{evaluation-fragmentation}}

No two dynamic weight decay papers share the same experimental
conditions. This fragmentation is compounded by publication bias against
null results (Lipton \& Steinhardt, 2019). We acknowledge our own
validation is currently limited in scale (CIFAR-only, ResNet-20);
large-scale replication is future work. Wang \& Aitchison (2024) showed
optimal WD scales as an EMA timescale \(\approx 1/\rho\) in our
notation---consistent with our conditional indistinguishability result.

\hypertarget{positioning-against-key-recent-works}{%
\subsubsection{2.4 Positioning Against Key Recent
Works}\label{positioning-against-key-recent-works}}

\textbf{vs.~Wang \& Aitchison (2024).} Their empirical scaling rule
focuses on WD \emph{magnitude}; our framework focuses on WD
\emph{modulation strategy} with regime-boundary theory and falsifiable
predictions.

\textbf{vs.~D'Angelo et al.~(2024).} Their finding is a static BN
reparameterization argument. Our Phi Invariance is a \emph{dynamic}
argument about AdamW's sign normalization. The NoBN ablation---not yet
completed---is needed to distinguish these mechanisms. We cannot
currently determine the primary driver of observed invariance.

\textbf{vs.~Chou (2025).} Their \(\lambda \propto \gamma\) schedule is a
specific temporal-axis modulation. Based on testing at a single
operating point (\(\rho = 0.5\)), it appears indistinguishable from
constant WD under AdamW. Regime II behavior remains untested.

\begin{center}\rule{0.5\linewidth}{0.5pt}\end{center}

\hypertarget{the-phi-modulator-framework}{%
\subsection{3. The Phi Modulator
Framework}\label{the-phi-modulator-framework}}

\hypertarget{formal-definition}{%
\subsubsection{3.1 Formal Definition}\label{formal-definition}}

We generalize AdamW by introducing the \textbf{Phi modulator}
\(\varphi\):
\[\boldsymbol{\theta}_{t+1} = \boldsymbol{\theta}_t - \eta_t \frac{\hat{\mathbf{m}}_t}{\sqrt{\hat{\mathbf{v}}_t} + \epsilon} - \lambda \cdot \varphi(t, \boldsymbol{\theta}_t, \mathbf{s}_t) \odot \boldsymbol{\theta}_t\]
where
\(\varphi : \mathbb{Z}_{\geq 0} \times \mathbb{R}^d \times \mathcal{S} \to \mathbb{R}^d_{\geq 0}\)
is a per-parameter, non-negative modulation function, \(\mathbf{s}_t\)
is the optimizer state (which may include \(\mathbf{g}_t\),
preconditioned updates \(\mathbf{u}_t\), or other statistics), and
\(\odot\) is element-wise multiplication.

\textbf{Budget reference convention:} \(\mathbb{E}[\varphi] = 1\) is the
normalization reference point. Methods with
\(\mathbb{E}[\varphi] \neq 1\) are valid modulators applying a different
total decay budget, quantified by BEM (Section 3.4). As a consequence,
no\_wd (\(\mathbb{E}[\varphi] = 0\)), half-lambda
(\(\mathbb{E}[\varphi] = 0.5\)), cosine
(\(\mathbb{E}[\varphi] \approx 0.4\)), and CWD
(\(\mathbb{E}[\varphi] \approx 0.5\)) all deviate from the reference.

\hypertarget{method-catalog}{%
\subsubsection{3.2 Method Catalog}\label{method-catalog}}

\textbf{Table 1: Method catalog.} The third argument \(\mathbf{s}_t\)
refers to optimizer state; for CWD this is specifically the
preconditioned update \(\mathbf{u}_t\). Methods with \(\dagger\) are
cataloged but not experimentally evaluated.

\begin{longtable}[]{@{}
  >{\raggedright\arraybackslash}p{(\columnwidth - 4\tabcolsep) * \real{0.1111}}
  >{\raggedright\arraybackslash}p{(\columnwidth - 4\tabcolsep) * \real{0.6528}}
  >{\raggedright\arraybackslash}p{(\columnwidth - 4\tabcolsep) * \real{0.2361}}@{}}
\toprule\noalign{}
\begin{minipage}[b]{\linewidth}\raggedright
Method
\end{minipage} & \begin{minipage}[b]{\linewidth}\raggedright
\(\varphi(t, \boldsymbol{\theta}, \mathbf{s})\)
\end{minipage} & \begin{minipage}[b]{\linewidth}\raggedright
Modulation Axis
\end{minipage} \\
\midrule\noalign{}
\endhead
\bottomrule\noalign{}
\endlastfoot
Constant (baseline) & \(\mathbf{1}\) & --- \\
CWD (hard) &
\(\mathbb{1}[\mathrm{sign}(\boldsymbol{\theta}) = \mathrm{sign}(\mathbf{u}_t)]\)
& Directional \\
SWD / AdamS & \(h(\|\mathbf{g}_t\|) \cdot \mathbf{1}\) &
Temporal-gradient \\
Cosine schedule & \(\tfrac{1}{2}(1 + \cos(\pi t / T)) \cdot \mathbf{1}\)
& Temporal \\
AdamWN\(^\dagger\) &
\(\max(0, 1 - \tau / \|\boldsymbol{\theta}\|) \cdot \mathbf{1}\) &
Target-norm \\
AlphaDecay\(^\dagger\) & \(\alpha_l \cdot \mathbf{1}_l\) (per layer
\(l\)) & Spatial \\
No-WD & \(\mathbf{0}\) & Ablation \\
Random mask & \(\mathrm{Bernoulli}(p) \cdot \mathbf{1}\) & Control \\
Half-lambda & \(0.5 \cdot \mathbf{1}\) & Budget control \\
\end{longtable}

AdamWN and AlphaDecay require architecture-specific hyperparameters;
evaluation deferred to future work. ADANA and AdamO jointly modify
weight decay and momentum/radial dynamics, placing them outside pure WD
modulation scope.

\hypertarget{budget-equivalence-normalization}{%
\subsubsection{3.3 Budget Equivalence
Normalization}\label{budget-equivalence-normalization}}

\textbf{Definition 1} (Budget Equivalence). \emph{Two strategies are
budget-equivalent if
\(\sum_{t} \lambda \cdot \mathbb{E}[\varphi_1(t, \boldsymbol{\theta}_t, \mathbf{s}_t)] = \sum_{t} \lambda \cdot \mathbb{E}[\varphi_2(t, \boldsymbol{\theta}_t, \mathbf{s}_t)]\).}

\hypertarget{diagnostic-metrics}{%
\subsubsection{3.4 Diagnostic Metrics}\label{diagnostic-metrics}}

\textbf{Budget Equivalence Metric (BEM):}
\[\mathrm{BEM}(\text{method}) = \frac{\bar{\lambda}_{\mathrm{eff}}^{\text{method}} - \bar{\lambda}_{\mathrm{eff}}^{\text{constant}}}{\bar{\lambda}_{\mathrm{eff}}^{\text{constant}}}\]
Verified values: half\_lambda \(= -0.500\), cosine\_schedule
\(\approx -0.600\), no\_wd \(= -1.000\).

\textbf{Coupling Stability Index (CSI):}
\[\mathrm{CSI}_{\mathrm{raw}} = \tfrac{1}{3} \cdot \mathrm{CV}(\|\mathbf{w}_t\|) + \tfrac{1}{3} \cdot \log(\kappa) + \tfrac{1}{3} \cdot \mathrm{CV}(\eta_{\mathrm{eff}})\]
Equal weights (\(1/3\) each) are a convention pending sensitivity
analysis. We report
\(\mathrm{CSI}_{\mathrm{rel}} = \mathrm{CSI}_{\mathrm{raw}} / \mathrm{CSI}_{\mathrm{constant}}\).
Note: \(\mathrm{CSI}_{\mathrm{rel}}\) is normalized to the constant
baseline on the same architecture; cross-architecture comparisons
require separate normalization.

\textbf{Alignment Informativeness Score (AIS):}
\[\mathrm{AIS} = \frac{1}{L}\sum_{l=1}^{L} \frac{H(\text{bin distribution of } |\cos(\boldsymbol{\theta}^{(l)}, \mathbf{g}^{(l)})|)}{\log(10)}\]
AIS \(\in [0, 1]\): values near 1 indicate high alignment diversity. AIS
is an intrinsic property of the network and loss landscape, not of the
weight decay method.

\begin{center}\rule{0.5\linewidth}{0.5pt}\end{center}

\hypertarget{theoretical-analysis-the-rho-regime-boundary}{%
\subsection{\texorpdfstring{4. Theoretical Analysis: The \(\rho\) Regime
Boundary}{4. Theoretical Analysis: The \textbackslash rho Regime Boundary}}\label{theoretical-analysis-the-rho-regime-boundary}}

\hypertarget{the-rho-lambdaeta-order-parameter}{%
\subsubsection{\texorpdfstring{4.1 The \(\rho = \lambda/\eta\) Order
Parameter}{4.1 The \textbackslash rho = \textbackslash lambda/\textbackslash eta Order Parameter}}\label{the-rho-lambdaeta-order-parameter}}

\[\rho = \frac{\bar{\lambda}}{\eta}\] where
\(\bar{\lambda} = \lambda \cdot \mathbb{E}[\varphi]\). At standard AdamW
settings, \(\rho = 0.5\).

\textbf{Remark 1} (Dual Characterization). \emph{The constraint radius
\(\tau^* = \eta/\lambda\) from Xie \& Li (2024) and the steady-state
ratio \(R_* \approx \lambda/\eta\) from Defazio (2025) are dual
characterizations of \(\rho\): \(\tau^* = 1/\rho\) and \(R_* = \rho\).
Neither prior work makes this connection explicit. Our contribution is
to operationalize \(\rho\) as an empirically testable regime boundary
parameter.} Wang \& Aitchison's (2024) EMA timescale \(\approx 1/\rho\),
adding a third perspective.

\hypertarget{the-phi-invariance-conjecture}{%
\subsubsection{4.2 The Phi Invariance
Conjecture}\label{the-phi-invariance-conjecture}}

\textbf{Conjecture 1} (Phi Invariance Trichotomy). \emph{There exist
\(0 < \rho_1 < \rho_2\) such that:}

\begin{itemize}
\tightlist
\item
  \textbf{Regime I} (\(\rho \leq \rho_1\)): \emph{Final loss difference
  for budget-equivalent schedules is
  \(O(\rho^2 \cdot V \cdot T \cdot \eta^2)\), implying accuracy
  differences \(< 0.1\%\) at \(\rho \approx 0.5\).}
\item
  \textbf{Regime II} (\(\rho_1 < \rho < \rho_2\)): \emph{Loss difference
  scales as
  \(O(\rho \cdot \sum_t |\lambda_t^{(1)} - \lambda_t^{(2)}| \cdot (1 - \delta_t))\);
  alignment-conditioned strategies can provide \(O(\rho)\) improvement.}
\item
  \textbf{Regime III} (\(\rho \geq \rho_2\)): \emph{All modulation
  strategies produce \(O(\rho)\) differences.}
\end{itemize}

Boundary estimates \(\rho_1 \approx 1\) and \(\rho_2 \approx 10\) are
order-of-magnitude placeholders to be determined by \(\lambda\)-sweep
experiments. \textbf{This conjecture is supported at a single operating
point (\(\rho = 0.5\), Regime I). Regimes II and III are untested
predictions.}

\textbf{Falsifiable predictions:} 1. At \(\rho = 0.5\): accuracy spread
\(< 0.5\%\). \emph{Observed}: \(0.25\%\) on CIFAR-10. Consistent. 2. At
\(\rho = 5\): accuracy spread \(1\)--\(3\%\). \emph{Untested}: requires
\(\lambda\)-sweep. 3. SGD should exhibit no Regime I. \emph{Observed}:
large effect sizes, consistent.

\textbf{Proof elements} (Appendix D, in preparation): The conjecture
rests on three lemmas. \emph{Lemma 1} (AdamW \(\ell_\infty\) norm
stability, from Xie \& Li 2024 Theorem 3.1). \emph{Lemma 2} (WD
perturbation bound via exponential damping). \emph{Lemma 3} (Regime I
negligibility: damped sum is
\(O(\rho^2 \cdot V \cdot T \cdot \eta^2)\)). Lemma 3 is the nontrivial
claim and is stated as an assumed bound pending formal derivation. The
critical assumption of Adam saturation is empirically verified for
\(> 80\%\) of parameters at epoch 100.

\hypertarget{why-sgd-differs}{%
\subsubsection{4.3 Why SGD Differs}\label{why-sgd-differs}}

Under AdamW, the preconditioned update approaches \(\pm 1\) per
coordinate, producing an implicit \(\ell_\infty\) constraint that
absorbs weight decay perturbations. SGD lacks this mechanism entirely.

\textbf{Disclosure on the 18.3\(\times\) ratio.}
\(\rho_{\text{AdamW}} = 0.5\) and
\(\rho_{\text{SGD}} = 0.05\)---10\(\times\) different. The observed
ratio reflects the combined effect of (1) the optimizer mechanism
difference and (2) the 10\(\times\) \(\rho\) difference. A
matched-\(\rho\) control is required to isolate the optimizer
contribution (Section 7.3, planned P1-c). We do not claim the ratio
measures a pure optimizer mechanism effect.

\textbf{Heuristic Observation 1} (SGD Alignment-Optimal Schedule,
conjectural). \emph{Under Sun et al.'s (CVPR 2025) fixed-\(\lambda\)
stability bound, extending to time-varying \(\lambda_t\) (formal gap
present), the optimal budget-allocated schedule is
\(\lambda_t^* \propto \delta_t / (1 - \delta_t)\).} Treated as a
heuristic motivation for CWD's binary mask approximation.

\begin{center}\rule{0.5\linewidth}{0.5pt}\end{center}

\hypertarget{experimental-setup}{%
\subsection{5. Experimental Setup}\label{experimental-setup}}

\hypertarget{implementation}{%
\subsubsection{5.1 Implementation}\label{implementation}}

All experiments use a unified \texttt{UnifiedAdamW} optimizer with
pluggable Phi modulator interface, ensuring identical optimizer
internals. Seven weight decay methods: \textbf{constant},
\textbf{cwd\_hard}, \textbf{swd}, \textbf{cosine\_schedule},
\textbf{random\_mask}, \textbf{half\_lambda}, \textbf{no\_wd}. Same
methods tested with SGD.

\textbf{Bug note.} BEM computation bug corrected (\texttt{HalfLambdaPhi}
reported \(0.000\) instead of \(-0.500\)). All results use corrected
implementation.

\hypertarget{training-configuration}{%
\subsubsection{5.2 Training
Configuration}\label{training-configuration}}

\textbf{Architecture.} ResNet-20 (\$\sim\$270K parameters, batch
normalization). VGG-16-BN pilot in Section 6.3.

\textbf{Optimizers.} AdamW: \(\eta = 10^{-3}\),
\(\lambda = 5 \times 10^{-4}\), \(\rho_{\text{AdamW}} = 0.5\). SGD:
\(\eta = 0.1\), \(\lambda = 5 \times 10^{-3}\),
\(\rho_{\text{SGD}} = 0.05\). Note: different \(\rho\) values by
standard configuration convention.

\textbf{Training.} 200 epochs, batch size 128, 3 seeds (42, 123, 456).
\textbf{Data integrity note}: CIFAR-100 SGD no\_wd has \(n = 1\) seed
(seed\_42 only); excluded from cross-dataset SGD statistics.

\textbf{Budget.} AdamW: 42 runs. SGD: 40 complete + 1 partial. VGG
pilot: 3 runs. Total: 87+.

\hypertarget{statistical-analysis-protocol}{%
\subsubsection{5.3 Statistical Analysis
Protocol}\label{statistical-analysis-protocol}}

\textbf{Pairwise comparisons.} Paired \(t\)-tests vs.~constant with Holm
correction. Cohen's \(d\) uses pooled standard deviation.

\textbf{Equivalence testing.} TOST with \(\delta = 0.5\%\). With
\(n = 3\): TOST power \(\approx 15\)--\(20\%\); MDE \(\approx 0.77\%\)
at 80\% power. We report non-significant differences \emph{consistent
with} equivalence---not proof. Formal TOST requires \(n \geq 5\).

\textbf{Bootstrap CIs.} 10,000-resample BCa 95\% CIs (approximate,
\(n = 3\)).

\textbf{Statistical honesty.} We distinguish: (i) significant results;
(ii) non-significant trends; (iii) observations consistent with but not
proving equivalence. All results reported regardless of direction.

\begin{center}\rule{0.5\linewidth}{0.5pt}\end{center}

\hypertarget{results}{%
\subsection{6. Results}\label{results}}

\hypertarget{adamw-conditional-indistinguishability}{%
\subsubsection{6.1 AdamW: Conditional
Indistinguishability}\label{adamw-conditional-indistinguishability}}

Figure 1 shows accuracy distributions under AdamW. \textbf{Table 2}
presents the main results. All findings are limited to ResNet-20 (batch
normalization, CIFAR scale) at \(\rho = 0.5\) (predicted Regime I; the
regime label is assumed---empirical boundaries require \(\lambda\)-sweep
validation).

\textbf{Table 2: AdamW results (ResNet-20, 200 epochs, 3 seeds).}

\begin{longtable}[]{@{}
  >{\raggedright\arraybackslash}p{(\columnwidth - 8\tabcolsep) * \real{0.1379}}
  >{\raggedright\arraybackslash}p{(\columnwidth - 8\tabcolsep) * \real{0.3276}}
  >{\raggedright\arraybackslash}p{(\columnwidth - 8\tabcolsep) * \real{0.3448}}
  >{\raggedright\arraybackslash}p{(\columnwidth - 8\tabcolsep) * \real{0.0690}}
  >{\raggedright\arraybackslash}p{(\columnwidth - 8\tabcolsep) * \real{0.1207}}@{}}
\toprule\noalign{}
\begin{minipage}[b]{\linewidth}\raggedright
Method
\end{minipage} & \begin{minipage}[b]{\linewidth}\raggedright
CIFAR-10 Acc. (\%)
\end{minipage} & \begin{minipage}[b]{\linewidth}\raggedright
CIFAR-100 Acc. (\%)
\end{minipage} & \begin{minipage}[b]{\linewidth}\raggedright
Cohen's \(d\) (C10)
\end{minipage} & \begin{minipage}[b]{\linewidth}\raggedright
BEM
\end{minipage} \\
\midrule\noalign{}
\endhead
\bottomrule\noalign{}
\endlastfoot
constant & 90.13 \(\pm\) 0.31 & 63.15 \(\pm\) 0.30 & --- & 0.000 \\
cosine\_schedule & 90.12 \(\pm\) 0.07 & 63.42 \(\pm\) 0.42 & 0.05 &
\(-0.600\) \\
cwd\_hard & 90.06 \(\pm\) 0.24 & 62.84 \(\pm\) 0.29 & 0.25 &
\(-0.490\) \\
random\_mask & 90.12 \(\pm\) 0.30 & 62.87 \(\pm\) 0.37 & 0.03 &
\(-0.500\) \\
half\_lambda & 90.09 \(\pm\) 0.28 & 62.91 \(\pm\) 0.47 & 0.14 &
\(-0.500\) \\
swd & 89.88 \(\pm\) 0.25 & 63.06 \(\pm\) 0.29 & 0.89 & varies* \\
no\_wd & 90.08 \(\pm\) 0.32 & 62.66 \(\pm\) 0.38 & 0.16 & \(-1.000\) \\
\end{longtable}

*SWD's BEM is time-varying; mean BEM \(\approx -0.15\) to \(-0.30\)
depending on gradient dynamics.

All pairwise \(p_{\text{adj}} > 0.05\) (Holm correction).

\textbf{Key observations:} 1. Accuracy range \(< 0.3\%\) on
CIFAR-10---within seed variance (\(\sigma \approx 0.25\)--\(0.30\%\)).
2. No significant pairwise differences (\(n = 3\), MDE
\(\approx 0.77\%\) at 80\% power). 3. No weight decay is statistically
indistinguishable from constant within our detection capability. 4.
Budget variation does not predict performance (Figure 3, BEM span
\([0.000, -1.000]\) with no accuracy correlation). Consistent with
Regime I at the single tested operating point. 5. Cosine schedule shows
notably lower variance on CIFAR-10 (\(\sigma = 0.07\%\)
vs.~\(\sim 0.30\%\)). This property does not replicate on CIFAR-100
(\(\sigma = 0.42\%\) for cosine); it should not be treated as a
consistent stability property without formal variance testing (Levene's
test is underpowered at \(n = 3\)).

\hypertarget{sgd-weight-decay-presence-matters}{%
\subsubsection{6.2 SGD: Weight Decay Presence
Matters}\label{sgd-weight-decay-presence-matters}}

Figure 2 presents the AdamW/SGD contrast. \textbf{Table 3} presents SGD
results on CIFAR-10.

\textbf{Table 3: SGD results (ResNet-20, CIFAR-10, 200 epochs, 3
seeds).}

\begin{longtable}[]{@{}llll@{}}
\toprule\noalign{}
Method & Accuracy (\%) & Cohen's \(d\) & \(p_{\text{adj}}\) (Holm) \\
\midrule\noalign{}
\endhead
\bottomrule\noalign{}
\endlastfoot
constant & 91.22 \(\pm\) 0.07 & --- & --- \\
cosine\_schedule & 91.20 \(\pm\) 0.12 & 0.17 & 0.869 \\
cwd\_hard & 90.87 \(\pm\) 0.43 & 1.08 & 0.349 \\
random\_mask & 90.77 \(\pm\) 0.45 & 1.37 & 0.218 \\
half\_lambda & 90.84 \(\pm\) 0.18 & 2.75 & 0.074 \\
swd & 90.71 \(\pm\) 0.19 & 3.48 & 0.054 \\
no\_wd & 90.30 \(\pm\) 0.10 & 10.29 & \textbf{0.002} \\
\end{longtable}

After Holm correction, \textbf{only constant vs.~no\_wd achieves
significance} (\(d = 10.29\), \(p_{\text{adj}} = 0.002\)). swd
(\(p_{\text{adj}} = 0.054\)) and half\_lambda
(\(p_{\text{adj}} = 0.074\)) do not reach \(\alpha = 0.05\), though
large effect sizes suggest the \(n = 3\) design lacks power.

\textbf{Key findings:} 1. Weight decay presence matters under SGD: 0.92
pp degradation without WD, \(d = 10.29\) (\(p_{\text{adj}} = 0.002\)).
The single most statistically reliable finding in this study. 2. Effect
sizes consistently large: five of six SGD comparisons yield \(|d| > 1\),
versus \(|d| < 1\) for all AdamW comparisons. 3. The
\textbf{18.3\(\times\) effect-size ratio} (SGD \(d = 10.29\) vs.~AdamW
\(d = 0.16\)). Bootstrap BCa 95\% CI: {[}\(12.1\times\),
\(28.7\times\){]} (approximate, \(n = 3\)). \textbf{This ratio conflates
two effects}: (1) optimizer mechanism difference and (2) 10\(\times\)
\(\rho\) difference (\(\rho_{\text{AdamW}} = 0.5\)
vs.~\(\rho_{\text{SGD}} = 0.05\)). A matched-\(\rho\) control is
required to isolate the optimizer contribution.

\hypertarget{cross-architecture-pilot-and-power-discussion}{%
\subsubsection{6.3 Cross-Architecture Pilot and Power
Discussion}\label{cross-architecture-pilot-and-power-discussion}}

VGG-16-BN pilot experiments (10 epochs, 1 seed, CIFAR-10) confirm
infrastructure readiness only.

\begin{longtable}[]{@{}llll@{}}
\toprule\noalign{}
Method & VGG-16-BN Acc. (10 ep) & BEM & CSI \\
\midrule\noalign{}
\endhead
\bottomrule\noalign{}
\endlastfoot
constant & 79.94\% & 0.000 & 0.996 \\
cwd\_hard & 80.30\% & \(-0.490\) & 1.011 \\
no\_wd & 80.61\% & \(-1.000\) & 0.988 \\
\end{longtable}

\textbf{Directional anomaly}: no\_wd (80.61\%) outperforms constant
(79.94\%) at 10 epochs---opposite to the predicted invariance direction.
This may reflect early-epoch dynamics or noise; conclusions require
200-epoch, 3-seed experiments.

\textbf{Statistical power assessment.} MDE \(\approx 0.77\%\) at 80\%
power. The CIFAR-100 AdamW spread of 0.76\% (cosine: 63.42\% vs.~no\_wd:
62.66\%) falls within our detection gap. \(n \geq 5\) seed extension
with TOST is required for the central indistinguishability claim.

\hypertarget{diagnostic-analysis}{%
\subsubsection{6.4 Diagnostic Analysis}\label{diagnostic-analysis}}

\textbf{BEM (Figure 3).} BEM spans \([0.000, -1.000]\) with no accuracy
correlation under AdamW---validating BEM as a descriptive
characterization tool consistent with the Regime I prediction at the
single tested operating point.

\textbf{Weight norm convergence} (Figure 5). Under AdamW, all seven
methods converge to weight norms in the range 95.89--97.04 (1.2\%
variation), despite \(10\times\) BEM variation. This supports---but does
not prove---the \(\ell_\infty\) constraint mechanism. The BN confound
remains: all architectures use batch normalization, and the NoBN
ablation is blocking for causal attribution (Section 7.3, Limitation 2).

\textbf{AIS.} Values range 0.25--0.50, indicating moderate alignment
diversity. The alignment signal CWD exploits carries some information
but is insufficient to produce detectable gains in the tested regime.

\begin{center}\rule{0.5\linewidth}{0.5pt}\end{center}

\hypertarget{discussion}{%
\subsection{7. Discussion}\label{discussion}}

\hypertarget{practical-implications}{%
\subsubsection{7.1 Practical
Implications}\label{practical-implications}}

Our findings support a practically actionable observation: \textbf{under
standard AdamW settings (\(\lambda = 5 \times 10^{-4}\),
\(\eta = 10^{-3}\), batch-normalized architectures at CIFAR scale), we
detected no statistically significant benefit from dynamic weight decay
strategies} (\(n = 3\), MDE \(\approx 0.77\%\); see Limitation 1). All
seven methods produce indistinguishable results within our detection
capability.

The finding that no weight decay is indistinguishable from constant WD
under AdamW at \(\rho = 0.5\) is noteworthy. Under our hypothesized
mechanism, AdamW's \(\ell_\infty\) constraint already provides norm
control. We caution: this rests on BN-equipped architectures; BN
scale-invariance (D'Angelo et al., 2024) may be an equally valid
explanation (Limitation 2).

\textbf{Boundary conditions.} Weight decay strategy matters when: (a)
using SGD (\textasciitilde18.3\(\times\) effect-size ratio, though
confounded by \(\rho\) difference; Limitation 6); (b) using high
\(\lambda\) (untested prediction); (c) BN-free architectures (open
question).

\textbf{Proposed decision heuristic (conjectured, single-point
validated):} 1. Compute \(\rho = \lambda/\eta\). 2. If \(\rho < 1\) with
AdamW and BN: constant WD likely sufficient. 3. If \(\rho > 1\) or using
SGD: dynamic strategies may help (untested regime prediction). 4. Use
BEM to verify differences are not budget effects.

\hypertarget{theoretical-implications}{%
\subsubsection{7.2 Theoretical
Implications}\label{theoretical-implications}}

\textbf{Dual characterization (Remark 1).} Xie \& Li's
\(\tau^* = 1/\rho\) and Defazio's \(R_* = \rho\) are dual
characterizations of the same quantity. Our contribution is to
operationalize \(\rho\) as a regime boundary parameter.

\textbf{Mechanistic hypothesis and its limits.} We hypothesize AdamW's
per-coordinate sign descent creates an implicit constraint set
determined by \(\rho\) alone, making budget-equivalent Phi modulations
produce indistinguishable trajectories. \textbf{Critical caveat:} all
experiments use BN architectures. We cannot distinguish: (a) AdamW's
\(\ell_\infty\) constraint absorbing WD perturbations; (b) BN
scale-invariance rendering WD irrelevant; (c) both jointly. The NoBN
ablation is designed to resolve this---and has not been run (Limitation
2). The mechanistic attribution should be read as a hypothesis.

\textbf{The no-WD equivalence finding.} That BEM \(= -1\) (no weight
decay) produces the same accuracy as constant WD challenges the
assumption that WD is an essential regularizer. Either the
\(\ell_\infty\) mechanism makes explicit WD redundant at \(\rho = 0.5\),
or BN scale-invariance makes WD irrelevant for BN-covered parameters.

\textbf{Consistency with Wang \& Aitchison (2024).} Their EMA timescale
\(\approx 1/\rho\) is consistent with our framework: if \(\rho\) remains
in Regime I across scales, the specific modulation is immaterial.

\textbf{Chou (2025) as a special case.} Based on our single tested
point, their \(\lambda \propto \gamma\) schedule is indistinguishable
from constant WD under AdamW in Regime I. Regime II behavior remains
untested.

\hypertarget{scope-limitations-and-future-experiments}{%
\subsubsection{7.3 Scope, Limitations, and Future
Experiments}\label{scope-limitations-and-future-experiments}}

\textbf{Verified scope.} CIFAR-10 and CIFAR-100 with ResNet-20 (270K
parameters, single architecture, BN), AdamW at \(\rho = 0.5\), SGD at
\(\rho = 0.05\), \(n = 3\) seeds.

\textbf{Limitations, ordered by priority:}

\begin{enumerate}
\def\labelenumi{\arabic{enumi}.}
\item
  \textbf{Statistical power (P0, blocking for equivalence claims).} MDE
  \(\approx 0.77\%\), TOST power \(\approx 15\)--\(20\%\). CIFAR-100
  spread of 0.76\% sits at detection edge. \emph{Planned}: \(n \geq 5\)
  seed extension with formal TOST.
\item
  \textbf{BN confound (P0, blocking for mechanistic claims).} All
  architectures use BN; cannot distinguish \(\ell_\infty\) mechanism
  from BN scale-invariance. \emph{Planned}: NoBN ablation
  (ResNet-20-NoBN, \textasciitilde1 GPU-hour).
\item
  \textbf{Single architecture (P0, blocking for scope claims).}
  ResNet-20 only. \emph{Planned}: VGG-16-BN full runs, larger models.
\item
  \textbf{Single \(\rho\) operating point (P0, blocking for regime
  theory).} Only \(\rho = 0.5\) tested; regime boundaries are untested
  predictions. \emph{Planned}: \(\lambda\) sweep
  (\(\rho \in \{0.05, 0.5, 5, 50\}\)).
\item
  \textbf{Theoretical incompleteness.} Appendix D (Lemma proofs) is in
  preparation; conjecture supported at one data point.
\item
  \textbf{\(\rho\) confound in 18.3\(\times\) ratio (P1).}
  \(\rho_{\text{SGD}} = 0.05\) and
  \(\rho_{\text{AdamW}} = 0.5\)---10\(\times\) different. Ratio
  conflates optimizer mechanism with \(\rho\) operating-point
  difference. \emph{Planned}: matched-\(\rho\) SGD control.
\item
  \textbf{Dataset scale (P2).} CIFAR only. \emph{Planned}: ImageNet
  pilot (ResNet-50).
\item
  \textbf{BEM bug (disclosed).} Corrected before all reported
  experiments.
\end{enumerate}

\textbf{Falsification criteria.} The Phi Invariance Conjecture would be
falsified if: at \(\rho = 0.5\) with AdamW and BN, any dynamic WD
strategy shows \(> 1\%\) accuracy improvement over constant at
\(n \geq 7\) seeds. A NoBN ablation where invariance breaks under AdamW
would challenge the \(\ell_\infty\) mechanism.

\begin{center}\rule{0.5\linewidth}{0.5pt}\end{center}

\hypertarget{conclusion}{%
\subsection{8. Conclusion}\label{conclusion}}

\hypertarget{summary-of-findings}{%
\subsubsection{8.1 Summary of Findings}\label{summary-of-findings}}

We introduced the Phi Modulator Framework, recovering seven major
dynamic weight decay methods as special cases along four modulation
axes. Across 7 methods, 2 datasets, and 2 optimizers on ResNet-20 (3
seeds each), we identified three central observations:

\textbf{First, no significant differences under AdamW}: at
\(\rho = 0.5\), all seven methods are statistically indistinguishable on
CIFAR-10 (range \(< 0.3\%\), all \(p_{\text{adj}} > 0.05\), \(n = 3\),
MDE \(\approx 0.77\%\)) and CIFAR-100. These observations are consistent
with conditional equivalence but do not constitute proof; formal TOST at
\(n \geq 5\) is required. Results are limited to batch-normalized
architectures at CIFAR scale (ResNet-20, 270K parameters).

\textbf{Second, optimizer-dependent sensitivity}: SGD yields constant
vs.~no\_wd Cohen's \(d = 10.29\) (\(p_{\text{adj}} = 0.002\))---an
effect-size ratio of approximately \(18.3\times\) (95\% CI:
{[}\(12\times\), \(29\times\){]}) relative to AdamW, reflecting the
combined effect of optimizer mechanism and the 100\(\times\) difference
in \(\rho\) values. We attribute the asymmetry to AdamW's
\(\ell_\infty\) constraint as a hypothesis; BN scale-invariance provides
an alternative mechanism pending NoBN ablation.

\textbf{Third, the \(\rho\) regime boundary}: the Phi Invariance
Conjecture, positing \(\rho\) as the order parameter governing WD
strategy importance. Supported only at \(\rho \approx 0.5\); regime
boundaries require \(\lambda\)-sweep validation. Connected to prior work
via the dual characterization in Remark 1.

Additionally, BEM, CSI, and AIS provide standardized diagnostic tools.
BEM tracked budget variation across \([0, -1]\) with no accuracy
correlation under AdamW, consistent with Regime I.

\hypertarget{future-work}{%
\subsubsection{8.2 Future Work}\label{future-work}}

Ordered by priority: (1) NoBN ablation (P0, blocking); (2) \(\rho\)
regime sweep (P0, central falsification); (3) \(n \geq 5\) seeds with
TOST (P1); (4) VGG-16-BN full runs (P1); (5) matched-\(\rho\) SGD
control (P1); (6) ImageNet pilot (P2); (7) formal proof of conjecture
(P3).

\hypertarget{broader-impact}{%
\subsubsection{8.3 Broader Impact}\label{broader-impact}}

Knowing \emph{when not to optimize} is as valuable as knowing \emph{how
to optimize}. For practitioners using AdamW at standard settings, our
results suggest that effort invested in dynamic weight decay strategy
selection is likely unnecessary at the scales and settings tested. The
\(\rho\)-based decision heuristic in Section 7.1 provides a starting
point for evaluating whether this finding applies to new settings.
Beyond weight decay, identifying regime boundaries via dimensionless
ratios may serve as a template for resolving similar proliferation
problems in optimizer design.

\begin{center}\rule{0.5\linewidth}{0.5pt}\end{center}

\hypertarget{references}{%
\subsection{References}\label{references}}

Chen, L. et al.~(2026a). Cautious Weight Decay. \emph{ICLR 2026}.

Chen, L. et al.~(2026b). AdamO: Decoupling Radial and Tangential Weight
Decay Dynamics. \emph{ICLR 2026}. {[}arXiv preprint, 2025{]}

Chou, S. (2025). Weight Decay Scheduling Proportional to Learning Rate.
\emph{arXiv preprint}.

D'Angelo, F. et al.~(2024). Why Weight Decay Is Never Useful as
Regularization in Modern Deep Learning. \emph{NeurIPS 2024}.

Defazio, A. (2025). Steady-State Analysis of Weight Decay in Adaptive
Optimizers. \emph{arXiv preprint}.

Ferbach, D. et al.~(2026). ADANA: Adaptive Norm and Weight Decay
Annealing. \emph{ICLR 2026}. {[}arXiv preprint, 2025{]}

Hanson, S. J. \& Pratt, L. Y. (1988). Comparing biases for minimal
network construction with back-propagation. \emph{NeurIPS 1988}.

He, K. et al.~(2016). Deep Residual Learning for Image Recognition.
\emph{CVPR 2016}.

He, Z. et al.~(2025). AlphaDecay: Module-Wise Weight Decay via Spectral
Density Analysis. \emph{arXiv preprint}.

Hoffer, E. et al.~(2018). Norm Matters: Efficient and Accurate
Normalization Schemes in Deep Networks. \emph{NeurIPS 2018}.

Ioannidis, J. P. A. (2005). Why Most Published Research Findings Are
False. \emph{PLoS Medicine}.

Kingma, D. P. \& Ba, J. (2015). Adam: A Method for Stochastic
Optimization. \emph{ICLR 2015}.

Kosson, A. et al.~(2023). Rotational Equilibrium: How Weight Decay
Balances Learning Across Neural Networks. \emph{arXiv preprint}.

Krizhevsky, A. (2009). Learning Multiple Layers of Features from Tiny
Images. \emph{Technical Report}.

Krogh, A. \& Hertz, J. (1991). A Simple Weight Decay Can Improve
Generalization. \emph{NeurIPS 1991}.

Lipton, Z. C. \& Steinhardt, J. (2019). Troubling Trends in Machine
Learning Scholarship. \emph{ACM Queue}.

Loshchilov, I. (2023). Weight Norm Control: Generalized Decoupled Weight
Decay. \emph{arXiv preprint}.

Loshchilov, I. \& Hutter, F. (2017). SGDR: Stochastic Gradient Descent
with Warm Restarts. \emph{ICLR 2017}.

Loshchilov, I. \& Hutter, F. (2019). Decoupled Weight Decay
Regularization. \emph{ICLR 2019}.

Sun, R. et al.~(2025). Alignment-Dependent Stability Bounds for SGD with
Weight Decay. \emph{CVPR 2025}.

Tian, Y. et al.~(2024). Selective Projection Decay for Fine-Tuning.
\emph{arXiv preprint}.

Wang, A. \& Aitchison, L. (2024). Optimal Weight Decay as EMA Timescale.
\emph{NeurIPS 2024}.

Xie, Z. et al.~(2023). Scheduled Weight Decay (AdamS). \emph{ICML 2023}.

Xie, Z. \& Li, Z. (2024). AdamW as \(\ell_\infty\)-Constrained
Optimization. \emph{NeurIPS 2024}.

\begin{center}\rule{0.5\linewidth}{0.5pt}\end{center}

\hypertarget{figures-and-tables}{%
\subsection{Figures and Tables}\label{figures-and-tables}}

\begin{itemize}
\tightlist
\item
  Figure 1: fig2\_accuracy\_comparison.png --- Accuracy distributions
  across 7 weight decay strategies under AdamW (bar chart with error
  bars, both datasets). Referenced in Sections 1.3 and 6.1.
\item
  Figure 2: fig2\_accuracy\_comparison.png --- AdamW vs.~SGD comparative
  panel showing the 18.3x effect-size contrast. Referenced in Sections
  1.3 and 6.2.
\item
  Figure 3: fig3\_bem\_vs\_accuracy.png --- BEM vs.~final accuracy
  scatterplot showing budget irrelevance under AdamW. Referenced in
  Sections 6.1 (Observation 4) and 6.4.
\item
  Figure 4: fig4\_diagnostic\_heatmap.png --- Diagnostic metric heatmap
  (BEM, CSI, AIS) across all methods. Referenced in Section 6.4.
\item
  Figure 5: fig5\_weight\_norm\_trajectories.png --- Weight norm
  trajectories under AdamW converging to common neighborhood. Referenced
  in Section 6.4.
\item
  Figure 6: fig6\_sgd\_vs\_adamw\_norms.png --- Weight norm trajectories
  contrasting SGD vs.~AdamW. Referenced in Section 6.4.
\item
  Table 1: inline --- Method catalog: Phi formulations for 9 weight
  decay strategies (Section 3.2).
\item
  Table 2: inline --- AdamW accuracy results, ResNet-20, 200 epochs, 3
  seeds (Section 6.1).
\item
  Table 3: inline --- SGD accuracy results, ResNet-20, CIFAR-10, 200
  epochs, 3 seeds (Section 6.2).
\end{itemize}

\begin{center}\rule{0.5\linewidth}{0.5pt}\end{center}

\emph{Integrated paper: Revision 1 (post-critique rounds 1--2)}
\emph{Date: 2026-03-18} \emph{Editor: Sibyl Editor Agent (sibyl-heavy)}

\end{document}
